\documentclass[profile=standard]{ipara-handbook}
\title{B06B｜期望、联合概率与边缘化：概率的积分语言}
\author{}
\date{}
\begin{document}
\maketitle

\begin{quote}
本册回答：\textbf{为什么概率论里到处都在积分，却每一次积分的“信息意义”并不相同？怎样把联合概率、边缘化、条件化、期望和卷积放进同一套累积语言，而又不把它们混成一件事？}
\end{quote}

\section{本册定位：积分是概率质量的汇总器，但汇总目标不同}
高数 Topic05/08 Own 一维与高维积分；概率 Topic04 Own 联合、边缘、条件和变量函数分布；概率 Topic05 Own 期望、方差等数字特征。B06B 只 Own “概率质量怎样调用积分完成不同信息操作”的跨科翻译。

最重要的四种语义是：
\[
\boxed{
\begin{array}{lll}
\text{区域概率} &:& \text{对目标区域内的质量求和};\\
\text{边缘化} &:& \text{沿不再保留的变量汇总};\\
\text{期望} &:& \text{按概率质量对数值加权};\\
\text{函数分布} &:& \text{沿同一输出值的原像/纤维汇总}. 
\end{array}}
\]

它们都使用积分，但对象、被积函数和信息损失完全不同。

\section{区域概率：最原始的“质量求和”}
若 $(X,Y)$ 有联合密度 $f_{X,Y}$，本册把“支撑” $S$ 理解为概率质量实际可能出现的取值区域；在普通连续题中，可先用 $f_{X,Y}(x,y)>0$ 的区域闭包来识别它。对目标事件区域 $D$，真正参与积分的是 $D\cap S$：
\[
P((X,Y)\in D)=\iint_{D\cap S} f_{X,Y}(x,y)\,\mathrm{d}x\,\mathrm{d}y.
\]
这时积分区域就是事件在取值空间中的几何表示与真实支撑的交。

所以概率积分第一原则不是“先定积分次序”，而是：
\[
\boxed{\text{先确定真实支撑 }S\text{ 与目标事件区域 }D\text{ 的交}.}
\]

积分限来自几何交集，不来自公式记忆。

\section{边缘化：积分在主动忘掉一个维度}
只关心 $X$ 时，对每个固定 $x$ 把所有可能的 $y$ 质量汇总：
\[
\boxed{f_X(x)=\int f_{X,Y}(x,y)\,\mathrm{d}y.}
\]

这不是“二维降成一维的计算技巧”，而是信息操作：
\[
(X,Y)\mapsto X.
\]
不同联合分布可能拥有同样的边缘，因此边缘化通常不可逆。依赖结构被丢失，不能由 $f_X,f_Y$ 单独重建联合分布。

这一点解释了为什么“边缘都一样”绝不等于“联合一样”或“变量独立”。

\section{条件化：不是忘掉，而是在一个截面上重新归一}
对连续型联合密度，若 $f_Y(y)>0$，
\[
\boxed{f_{X\mid Y}(x\mid y)=\frac{f_{X,Y}(x,y)}{f_Y(y)}.}
\]

边缘化是把 $Y$ 忘掉；条件化恰恰相反，是把 $Y=y$ 的信息加入当前世界，再对相应截面重新归一。

因此必须主动阻断：
\[
\boxed{\text{Marginalize}\neq\text{Condition}.}
\]
虽然两者都从联合密度出发，但信息方向相反。

\section{期望：积分在做“值 × 质量”的加权汇总}
若目标不是完整分布，只问函数 $g(X,Y)$ 的平均值，则
\[
\boxed{E[g(X,Y)]=\iint g(x,y)f_{X,Y}(x,y)\,\mathrm{d}x\,\mathrm{d}y}
\]
（在相应可积条件下）。

这说明期望的积分不是在“求某个区域的概率”，而是在每个位置先把数值 $g(x,y)$ 乘以当地概率质量，再汇总。

所以只问 $E[g(X)]$ 时，常常无需先求 $Y=g(X)$ 的完整分布；这不是捷径，而是目标只需要一个线性统计泛函。

\section{Fubini/Tonelli 背景：为什么可以逐层汇总}
二维积分之所以能写成“先一维、再另一维”，依赖积分理论的逐层汇总原则。对非负函数可由 Tonelli 保证交换/逐层积分的合法性；对绝对可积函数可由 Fubini 处理。

数学一不要求掌握测度论版本，但需要保留一个边界意识：\textbf{积分次序可以换，不是因为符号长得对称，而是因为可积性与区域表示满足相应条件。}

这也解释概率中的全期望：
\[
E[X]=E(E[X\mid Y])
\]
本质上是把同一联合质量按层先汇总、再把各层重新合成整体。

\section{卷积：沿“同一输出值的纤维”汇总}
令
\[
Z=X+Y.
\]
要得到 $Z=z$ 附近的概率质量，必须把所有满足
\[
x+y=z
\]
的联合输入贡献加起来。

若 $X,Y$ 独立且有密度，则
\[
\boxed{f_Z(z)=\int f_X(x)f_Y(z-x)\,\mathrm{d}x.}
\]
卷积并不是一个与重积分无关的新公式；它是在约束纤维 $x+y=z$ 上汇总概率质量后得到的一维表达。

没有独立性时，正确起点仍是联合密度：
\[
f_Z(z)=\int f_{X,Y}(x,z-x)\,\mathrm{d}x
\]
（结合真实支撑理解）。独立性只允许把联合密度分解成两个边缘密度的乘积。

\section{母例：三角支撑上一次看清区域、边缘与期望}
设联合密度
\[
f_{X,Y}(x,y)=2,
\qquad 0<y<x<1.
\]

\subsection{边缘化}
固定 $x$，允许的 $y$ 是 $0<y<x$，所以
\[
f_X(x)=\int_0^x2\,\mathrm{d}y=2x,
\qquad 0<x<1.
\]

\subsection{期望}
若只问 $E[X]$，可直接对联合质量加权：
\[
E[X]=\int_0^1\int_0^x 2x\,\mathrm{d}y\,\mathrm{d}x
=\int_0^1 2x^2\,\mathrm{d}x=\frac23.
\]
也可以先求 $f_X$ 再积分。两条路线结果相同，但信息路径不同。

这个例子说明：同一张支撑图可以服务多个概率操作，但每次必须明确“保留什么、消去什么、加权什么”。

\section{B06B 串起哪些章节}
\begin{tabular}{p{0.23\linewidth}p{0.30\linewidth}p{0.39\linewidth}}
\hline
Owner / 下游 & 提供对象 & B06B 的翻译责任\\
\hline
高数 Topic05/08 & 一维/高维累积、换序 & 把积分解释成概率质量汇总，并保留可积性/区域边界\\
概率 Topic04 & 联合、边缘、条件、卷积 & 区分“忘掉信息、加入信息、沿纤维汇总”\\
概率 Topic05 & 期望、方差、全期望 & 把数字特征解释成对完整分布的有损加权摘要\\
B07 & 随机变量变换 & 提供“输出值对应一组原像，质量需在原像上汇总”的背景\\
\hline
\end{tabular}

\section{反桥梁：积分符号相同不代表对象相同}
\begin{tabular}{p{0.20\linewidth}p{0.04\linewidth}p{0.20\linewidth}p{0.48\linewidth}}
\hline
A & & B & 真正区别\\
\hline
边缘化 & $\neq$ & 条件化 & 前者忘掉变量，后者固定信息后重新归一\\
区域概率 & $\neq$ & 期望 & 前者被积函数通常是密度；后者还乘目标数值函数\\
卷积 & $\neq$ & 密度直接相加 & 卷积沿同一和值的原像汇总贡献\\
边缘相同 & $\neq$ & 联合相同 & 边缘化丢失依赖结构，一般不可逆\\
\hline
\end{tabular}

\section{调用信号与第一观察}
看到联合密度题时先问四件事：
\begin{enumerate}
\item 目标是概率、边缘、条件、期望，还是新函数的分布？
\item 当前支撑是什么，固定一个变量后另一变量允许在哪个截面？
\item 这次积分是在“消去变量”“按值加权”还是“沿原像纤维汇总”？
\item 最终结果应满足什么独立校验：归一化、概率范围、期望量纲、边缘再积分为 $1$？
\end{enumerate}

\section{一页压缩}
B06B 的母句不是“概率题要会积分”，而是：
\[
\boxed{\text{概率积分 = 对概率质量做有明确语义的汇总。}}
\]

做题前必须给积分贴标签：
\[
\boxed{\text{区域质量 / 消元 / 条件归一 / 加权摘要 / 原像汇总}.}
\]
标签一旦明确，积分限、被积对象与信息损失才不会混乱。

\end{document}
